\chapter{Measuring the Simulation}

The last act measures a measurement. The preceding chapter simulated an experiment that cannot take place in
that exact form: an electron, a second electron, a Cooper-pair boundary, a carried corridor residue, a decider,
and an adjacent charge aperture are left inside the Lean apparatus and advanced as a finite study. The result is
an internal alpha receipt. It expresses \(\alpha\) as a combination of ratios:
\[
  \alpha = \frac{e^2}{4\pi},
\]
where the numerator is selected from a typed neighborhood of rational readings of charge square. The denominator
is the finite \(4\pi\) rotation scale chosen for the apparatus, and the selected face is supplied by the carried
state of the corridor.

That is only half the story. A simulation also has a cost. If the apparatus asks Lean to elaborate the object
that represents a reading, Lean spends heartbeats. Those heartbeats are not a theorem about nature. They are the
instrument reading of the formal machine that produced the receipt. The experiment is therefore measured twice:
first as the alpha-like number produced by the simulated corridor, and second as the effort spent by the
elaborator to make that corridor into a checked object.

The companion file for this act is
\[
  \texttt{device/Measurement/Meanwhile51.lean}.
\]
It imports the finished alpha receipt, gathers the earlier calibration receipts, and prints one final object:
\begin{quote}\small
\begin{verbatim}
#eval simulationMeasurementReport
\end{verbatim}
\end{quote}
This is the closing ledger for Volume One. It does not invent a new physics step. It asks whether the simulated
experiment has itself been measured and calibrated.

\section{The Instrument}

The instrument is Lean's elaborator heartbeat counter. The measured operation is not ordinary execution time.
Wall-clock time depends on the machine, load, scheduler, cache state, and many details outside the formal object.
Heartbeats are closer to what this apparatus needs: a bounded counter of elaborator work. They are still
implementation readings, and another Lean version or another elaborator may give different numbers, but they
are reproducible enough to calibrate inside one build.

The measurement command was introduced earlier:
\begin{quote}\small
\begin{verbatim}
register_heart_rate_as name term
\end{verbatim}
\end{quote}
Operationally, the command elaborates a quoted term, reduces it to weak-head normal form, reads the heartbeat
counter around that work, and then registers the resulting natural number as a Lean definition. In the language
of the book, it turns effort into a typed number.

The protocol is deliberately plain:
\begin{enumerate}
\item perform a warmup read;
\item discard the warmup;
\item take a second read as the measured value;
\item take a third read as the stability check;
\item accept the calibration when the second and third reads agree within one heartbeat.
\end{enumerate}
The final report records the same rule in words: discard the warmup, take the second read, and check the third
read within one heartbeat.

\section{The Heartbeat Calibration}

The first calibration is the control term. The report reads:
\begin{quote}\small
\begin{verbatim}
control := { warmup := 1618,
             read2 := 1602,
             read3 := 1602,
             stable := true }
\end{verbatim}
\end{quote}
The warmup spends \(1618\) heartbeats. The two measured reads both spend \(1602\). This is the baseline: the
instrument can repeat a small term without drift after the warmup has settled.

The next reads decompose the apparatus into four stages:
\begin{center}\small
\begin{tabular}{c|r|r|r|c}
operation & warmup & read2 & read3 & stable\\
\hline
base & 309 & 309 & 309 & true\\
projection & 4490 & 4462 & 4462 & true\\
recovered & 2164 & 2164 & 2164 & true\\
decision & 1376 & 1376 & 1376 & true
\end{tabular}
\end{center}
These are not decorative costs. They are the pieces from which the book's effort arithmetic is made. The base
pulse is \(309\). Projection contributes a large positive delta:
\[
  4462 - 309 = 4153.
\]
Recovery is cheaper than projection:
\[
  2164 - 4462 = -2298.
\]
Decision then adds:
\[
  1376 - 309 = 1067.
\]
The report records the predicted heart rate as
\[
  3231.
\]
The important point is that every number in this paragraph is a measured elaborator number, not a guessed
coefficient.

The target terms close the calibration:
\begin{center}\small
\begin{tabular}{c|r|r|r|c}
target & warmup & read2 & read3 & stable\\
\hline
\(\texttt{B2}\) & 2165 & 2165 & 2166 & true\\
\(\texttt{driver}\) & 2164 & 2165 & 2164 & true\\
\(\texttt{driverDef}\) & 1603 & 1603 & 1602 & true
\end{tabular}
\end{center}
The corresponding binding energies are:
\[
  -1066,\qquad -1066,\qquad -1628.
\]
In ordinary prose: the target can cost less than the separated stage accounting. The apparatus keeps that
difference as binding energy. It is the elaborator analogue of saying that a composed object may be cheaper to
hold than the sum of the pieces used to describe it.

\section{Charge as Elaboration Cost}

The electron charge receipt uses the same idea one level higher. The apparatus does not merely assign a charge
sign by convention. It measures the cost of elaborating the electron reading and then orients that cost by the
electron convention:
\begin{quote}\small
\begin{verbatim}
elaborationCost := 2337
chargeMagnitude := 2337
signedCharge := -2337
chargeUnit := -1
chargeOrientation := electron
costIsChargeMagnitude := true
chargeUnitMatchesElectronConvention := true
\end{verbatim}
\end{quote}
This is the place where the phrase ``the elaboration cost is the charge of the electron'' becomes a receipt.
The measured magnitude is \(2337\). The electron convention supplies the sign, so the signed charge is
\(-2337\). The check fields say that the magnitude and convention agree with the object being read.

This is not the elementary charge in coulombs. It is a formal charge reading in the internal apparatus. The
physical route enters only when a later experiment calibrates that internal reading against laboratory scales.
For Volume One, the measured fact is narrower and sharper: the electron object has a stable elaborator cost, and
the apparatus can interpret that cost as its internal charge magnitude.

\section{Machine Epsilon}

The machine also needs a smallest visible change. That reading is supplied by the electron machine epsilon
report. It compares a composed trace route with a direct trace route. The two routes produce the same trace
number:
\begin{quote}\small
\begin{verbatim}
composedTraceNumberDepth := 2
directTraceNumberDepth := 2
traceNumbersEqual := true
\end{verbatim}
\end{quote}
but they do not spend the same effort:
\begin{center}\small
\begin{tabular}{c|r|r|r|r|c}
route & warmup & cost & stability read & settling width & within epsilon\\
\hline
composed & 10103 & 10087 & 10083 & 4 & true\\
direct & 7521 & 7519 & 7518 & 1 & true
\end{tabular}
\end{center}
The composed route is larger by
\[
  10087 - 7519 = 2568.
\]
The apparatus records that this delta resolves above machine epsilon:
\begin{quote}\small
\begin{verbatim}
deltaResolvesAboveMachineEpsilon := true
machineEpsilonHeartbeat := 4
machineEpsilonScale := 10087
machineEpsilonScaledAt18 := 396550014870625
\end{verbatim}
\end{quote}
In decimal form, the scaled epsilon is approximately
\[
  0.000396550014870625.
\]
That number is not a universal constant. It is the smallest settled heartbeat change seen by this measurement
route, normalized by the composed cost scale. It says how fine the elaborator instrument is in this local
configuration.

\section{The Orbit Effort Curve}

The simulated experiment advances an orbit-carry state. The heartbeat cost rises as the state is carried farther
around the loop:
\begin{center}\small
\begin{tabular}{r|r|c|r}
steps & heartbeats & selected face & residue numerator\\
\hline
0 & 801663 & charge & 0\\
1 & 842925 & mass & 111111111110\\
2 & 863364 & value & 222222222219\\
4 & 879432 & value & 111109777771\\
8 & 918377 & mass & 333332444433\\
16 & 983774 & mass & 333330222199\\
32 & 1111099 & value & 333329777733\\
64 & 1372453 & value & 111111111034\\
128 & 1889780 & value & 444289
\end{tabular}
\end{center}
This is the stress-strain curve of the formal orbit. The independent variable is how far the carried state has
been advanced. The dependent variable is the elaborator effort required to report that state. The face and
residue columns say where the corridor landed.

The two-orbit slip grid supplies a smaller local measurement:
\begin{quote}\small
\begin{verbatim}
twoOrbitSlipGridHeartbeat := 921040
twoOrbitSlipGridBracketCount := 2
\end{verbatim}
\end{quote}
The grid has six cells and brackets two slip points. The first bracket lies between steps \(3\) and \(4\). The
second lies between steps \(5\) and \(6\). The brackets are not real-number limits imposed from outside. They
are partition labels produced by the apparatus and then interpreted as rational neighborhoods.

\section{Alpha as a Combination of Ratio Types}

The final alpha reading is not one ratio. It is a nested interpretation of several ratios.

First, there is the charge-square ratio:
\[
  \alpha_i = \frac{e_i^2}{4\pi}.
\]
The numerator \(e_i^2\) is selected from the lower, target, or upper aperture reading. The selected face is a
typed face of the corridor: charge, mass, or value.

Second, there is the rotation ratio. The denominator \(4\pi\) is represented by a finite scaled integer. The
point is not to possess the real number \(\pi\) inside Lean. The point is to choose a finite rotation numerator,
run the apparatus against that finite scale, and record exactly which finite scale was used.

Third, there is the sampling ratio. In the \(10000\), \(20000\), and \(40000\) trial studies, the decider samples
the three faces with these counts:
\begin{center}\small
\begin{tabular}{r|r|r|r|r}
trials & charge & mass & value & mean \(\alpha\) scaled at \(10^{18}\)\\
\hline
10000 & 3389 & 3337 & 3274 & 7298680426529076\\
20000 & 6621 & 6655 & 6724 & 7298680845095521\\
40000 & 13349 & 13332 & 13319 & 7298680692889066
\end{tabular}
\end{center}
These counts are close to balanced, but not forced to be exactly balanced. The mean moves because the carried
state, residue, and decider all move.

Fourth, there is the refinement ratio. The Richardson report uses a refinement ratio of \(2\) and first-order
extrapolation. The two signed slips are
\[
  +418566445,\qquad -152206455
\]
at scale \(10^{18}\). Since the signs cross, the two Richardson estimates bracket the extrapolated value:
\[
  7298680540682611
  \le
  \alpha_{\mathrm{Rich}}
  \le
  7298681263661966.
\]
The midpoint is
\[
  7298680902172288
\]
at scale \(10^{18}\), or approximately
\[
  0.007298680902172288.
\]
The report also gives the source comparison:
\[
  \texttt{estimateSourceAbsDiffScaledAt18} = 12332646515.
\]
In decimal form that is about \(0.000000012332646515\).

Fifth, there is the effort ratio. The alpha reading did not arrive for free. It arrived through a calibrated
elaborator whose local machine epsilon is approximately
\[
  \frac{4}{10087}.
\]
So the final act measures not only which alpha-like number the simulated apparatus reads, but how costly the
apparatus was when it produced that reading.

\section{How to Reproduce the Last Act}

Run:
\begin{quote}\small
\begin{verbatim}
cd device
lake build Measurement.Meanwhile51
\end{verbatim}
\end{quote}
The final line of interest is:
\begin{quote}\small
\begin{verbatim}
info: Measurement/Meanwhile51.lean:...:
{ name := "measuring-the-simulation",
  heartbeatCalibration := ...,
  electronElaborationCharge := ...,
  electronMachineEpsilon := ...,
  orbitCarryEffort := ...,
  twoOrbitSlipGridHeartbeat := 921040,
  twoOrbitSlipGridBracketCount := 2,
  alphaRichardson := some ... }
\end{verbatim}
\end{quote}
If imported files print their own receipts first, keep scrolling until the final report named
``measuring-the-simulation'' appears.

To calibrate another Lean build, do not copy the numbers from this chapter. Copy the protocol. Record the
warmup, second read, and third read for each term. Check stability. Record the local machine epsilon. Record the
orbit effort curve. Then run the alpha study and Richardson report. Two systems can be compared only after both
have reported their own effort scale.

This is why different programming languages and elaborators matter. A Haskell prototype, a Rust evaluator, a
Coq development, an Agda development, and a Lean development may all implement the same corridor arithmetic, but
their effort instruments are different. The numerical alpha receipt belongs to the simulated apparatus. The
heartbeat receipt belongs to Lean. A serious comparison must keep those two layers separate and then calibrate
the map between them.

\section{What Has Been Measured}

The book has not measured the laboratory fine-structure constant. It has measured the simulation that the book
constructed. More precisely, it has measured:
\begin{enumerate}
\item the finite charge-square aperture used to recover an internal \(\alpha\);
\item the carried decider distribution over charge, mass, and value faces;
\item the Richardson bracket produced by three finite studies and two signed slips;
\item the elaborator effort required to construct representative readings;
\item the local machine epsilon of the elaborator measurement route;
\item the orbit effort curve showing how cost rises as the state is carried.
\end{enumerate}
That is the promised calibration. The simulated experiment is now an object with a reading and a cost. The reader
can rerun it, change the scales, change the trial count, change the route, or replace the internal ratios with
laboratory ledgers. The apparatus will expose the consequence as a new typed receipt.

Volume One therefore ends with the machine in a measurable state. The alpha-like number is on the page. The
ratios that produced it are on the page. The effort spent by the elaborator is on the page. The calibration rule
is simple enough to repeat: discard the warmup, take the second read, verify the third read, and report the local
epsilon. Nothing is hidden behind an appeal to mystery. The simulation was measured by the same discipline that
constructed it.
